\section{Genuine Markoff sources and their balanced cores}\label{sec:markoff-sources}

The reflection machine will accept continued fractions from a larger
language than the one supplied by Markoff occurrences. We now identify
the genuine sources inside that language. The distinction matters:
balance controls long factors, whereas the exact occurrence also fixes
the initial phase and the terminal digit.

\subsection{The mechanical word and the Cohn normalization}

Fix relatively prime integers $0<P<Q$. The lower Christoffel word
$C_{P/Q}=c_0\cdots c_{Q-1}$ is defined by
\begin{equation}\label{eq:source-mechanical}
 c_i=b\quad\Longleftrightarrow\quad
 \left\lfloor\frac{(i+1)P}{Q}\right\rfloor
 -\left\lfloor\frac{iP}{Q}\right\rfloor=1.
\end{equation}
It has $P$ letters $b$ and $Q-P$ letters $a$; thus $Q$ is its
total length. Its first letter is $a$ and its last is $b$, so write
$C_{P/Q}=azb$. The interior $z$ is a palindrome. Indeed, for
$1\leqslant j\leqslant Q-1$, coprimality gives
$$
 \left\lfloor\frac{(Q-j)P}{Q}\right\rfloor
 =P-1-\left\lfloor\frac{jP}{Q}\right\rfloor.
$$
Applying this identity to two consecutive interior indices shows that
$c_i=c_{Q-1-i}$ for $1\leqslant i\leqslant Q-2$.
Such interiors are called central words; see
\cite{BertheDeLucaReutenauer2008} for their standard word-theoretic
description.

Put
\begin{equation}\label{eq:source-cohn-matrices}
 A=M(1)^2=\begin{pmatrix}2&1\\1&1\end{pmatrix},\qquad
 B=M(2)^2=\begin{pmatrix}5&2\\2&1\end{pmatrix},
\end{equation}
and let $\mu$ be the word morphism with $\mu(a)=A$ and $\mu(b)=B$.
The Cohn--Christoffel correspondence identifies the normalized
Markoff occurrence at $P/Q$ with the label
\begin{equation}\label{eq:source-label}
 m_{P/Q}=\frac13\operatorname{tr}\mu(C_{P/Q})
        =\mu(C_{P/Q})_{12}.
\end{equation}
This is the classical correspondence with Markoff triples, not an
assertion that distinct fractions have distinct labels
\cite{Reutenauer2009}.
The two initial Markoff values $1$ and $2$ lie outside this proper-word
normalization and will not be needed below.

\subsection{The primitive column and its exact cost}

Use the digit product $M(q)$ from \eqref{eq:digit-matrix}. Identify
$a,b$ with the numerical input letters $1,2$, so that $\chi(a)=11$ and
$\chi(b)=22$. Reversal is denoted by a tilde.

\begin{proposition}[The genuine source dictionary]\label{prop:genuine-source-dictionary}
For every reduced $0<P<Q$, there is a canonical continued-fraction
word
\begin{equation}\label{eq:source-half}
 q=\varepsilon\chi(w)2,\qquad
 \varepsilon\in\{\varnothing,1,2\},
\end{equation}
whose first column $v=(y,x)^{\mathsf T}$ satisfies
$y>x>0$, $\gcd(x,y)=1$, and $x^2+y^2=m_{P/Q}$.
The core $w$ is a finite balanced binary word. Moreover,
\begin{equation}\label{eq:source-counts}
 \begin{aligned}
 P&=|q|_2=1+\mathbf1_{\{\varepsilon=2\}}+2|w|_b,\\
 Q&=|q|+1=2+\mathbf1_{\{\varepsilon\ne\varnothing\}}+2|w|.
 \end{aligned}
\end{equation}
and its subtractive Euclidean cost is
\begin{equation}\label{eq:source-cost}
 \ell(v)=\sum_{j=1}^s q_j-1=P+Q-2.
\end{equation}
The complete word is determined by the central palindrome through
\begin{equation}\label{eq:source-full-palindrome}
 \widetilde q\,q=2\chi(z)2.
\end{equation}
\end{proposition}

\begin{proof}
Split the central palindrome at its middle, in exactly one of the forms
$$
 z=\widetilde w\,w,\qquad
 z=\widetilde w\,a\,w,\qquad
 z=\widetilde w\,b\,w.
$$
The corresponding choices are $\varepsilon=\varnothing,1,2$.
Since $A$ and $B$ are symmetric, in all three cases
\begin{equation}\label{eq:source-half-factorization}
 \mu(z)=H^{\mathsf T}H,\qquad
 H=\begin{cases}
       \mu(w),&\varepsilon=\varnothing,\\
       M(\varepsilon)\mu(w),&\varepsilon\in\{1,2\}.
    \end{cases}
\end{equation}
Let $e=(2,1)^{\mathsf T}=M(2)(1,0)^{\mathsf T}$. The first row of
$A$ is $e^{\mathsf T}$ and the second column of $B$ is $e$.
Consequently
$$
 \mu(azb)_{12}=e^{\mathsf T}\mu(z)e=\|He\|^2.
$$
For completeness, if
$\mu(z)=\begin{pmatrix}a&b\\b&d\end{pmatrix}$, direct multiplication
gives
$$
 \mu(azb)_{12}=4a+4b+d,
 \qquad \operatorname{tr}\mu(azb)=3(4a+4b+d),
$$
which also verifies the equality of the two label conventions in
\eqref{eq:source-label}.

The matrix $H$ is integral and unimodular. Therefore $v=He$ is
primitive. Its coordinates are positive, with the first larger than
the second: every digit matrix $M(j)$ sends a positive column to one
with this ordering. By construction $HM(2)=M(q)$, proving the norm
claim. The last digit of $q$ is $2$, so $q$ is the canonical finite
continued fraction of $y/x$. This also includes the empty center,
where $q=2$ and $v=(2,1)^{\mathsf T}$.

Reversing the digit word in each of the three middle cases gives
\eqref{eq:source-full-palindrome} literally. In particular, the full
word has length $2+2|z|=2Q-2$ and contains
$2+2|z|_b=2P$ digits equal to $2$. Its two reversed halves have equal
length and equal digit counts. Hence $|q|=Q-1$ and $|q|_2=P$,
which gives \eqref{eq:source-counts}.

Equation \eqref{eq:euclidean-cost} gives $\ell(v)=\sum q_j-1$.
Since all digits are $1$ or $2$, their sum is $|q|+|q|_2=Q-1+P$,
proving \eqref{eq:source-cost}.

Finally, a binary word is \emph{balanced} if any two of its contiguous
factors of equal length have numbers of $b$'s differing by at most one.
For a length-$r$ factor of $C_{P/Q}$ starting at $i$, the number of
$b$'s is, by telescoping \eqref{eq:source-mechanical},
$$
 \left\lfloor\frac{(i+r)P}{Q}\right\rfloor
 -\left\lfloor\frac{iP}{Q}\right\rfloor.
$$
For fixed $r$ this is either $\lfloor rP/Q\rfloor$ or
$\lceil rP/Q\rceil$. Thus $C_{P/Q}$ is balanced, as is every
contiguous factor. The right half $w$ of $z$ is one such factor.
\end{proof}

The proof explains both roles of the source grammar. Pairing the digits
produces an integer Gram square and an exact cost; the mechanical word
supplies balance. No cyclic balance of an arbitrary central half is
asserted. In particular, a repeated segment of a finite core must be
analyzed with its actual initial and terminal states.

Conversely, the conditions that $q$ has the form
\eqref{eq:source-half}, that $w$ is balanced, and that the counts in
\eqref{eq:source-counts} are coprime are only necessary conditions in
this paper. An exact recognition test must also compare
$\widetilde q\,q$ with $2\chi(z_{P/Q})2$, where $z_{P/Q}$ is the
interior computed from \eqref{eq:source-mechanical}. If this equality
holds, the first-column construction proves that the source is exactly
the occurrence at $P/Q$. The reflection arguments below apply to a
larger language when this final recognition test is not imposed.

\subsection{How thin the genuine language is}

Balance alone is a weak constraint compared with being a genuine core. The
classical enumeration of the finite factors of Sturmian words makes the
comparison exact. Those factors are precisely the balanced words, and there
are $1+\sum_{i=1}^{n}(n-i+1)\varphi(i)$ of length $n$, with $\varphi$ the
Euler function; Mignosi proved this, and Berstel and Pocchiola gave a
geometric proof by duality and Euler's relation for planar graphs
\cite{Mignosi1991,BerstelPocchiola1993}.

\begin{proposition}[The genuine cores of a given length]
\label{prop:thin-language}
Fix $n\geqslant1$. The pairs $(\varepsilon,w)$ produced by
\cref{prop:genuine-source-dictionary} with $\lvert w\rvert=n$ are in
bijection with the reduced fractions $P/Q$, $0<P<Q$, having
$Q\in\{2n+2,\,2n+3\}$. There are therefore exactly
\begin{equation}
 \varphi(2n+2)+\varphi(2n+3)\leqslant 3n+3
 \label{eq:genuine-core-count}
\end{equation}
of them, whereas the balanced words of length $n$ number
$1+\sum_{i=1}^{n}(n-i+1)\varphi(i)$, which is asymptotic to $n^3/\pi^2$.
\end{proposition}

\begin{proof}
By \eqref{eq:source-counts} one has
$Q=2+\mathbf1_{\{\varepsilon\ne\varnothing\}}+2n$, so $\lvert w\rvert=n$
forces $Q=2n+2$ when $\varepsilon=\varnothing$ and $Q=2n+3$ otherwise;
conversely both values occur, the parity of $Q$ being the parity of
$\lvert z\rvert=Q-2$ that decides the three middle cases in the proof of
\cref{prop:genuine-source-dictionary}. That proof assigns to each reduced
$P/Q$ exactly one pair, so the map is well defined. It is injective because
the pair recovers the central word, $z=\widetilde w\,w$,
$\widetilde w\,a\,w$ or $\widetilde w\,b\,w$ according as $\varepsilon$ is
$\varnothing$, $1$ or $2$, and distinct reduced fractions have distinct
central words. Counting reduced fractions with a prescribed denominator
gives \eqref{eq:genuine-core-count}; the bound follows from
$\varphi(2n+2)\leqslant n+1$, valid because $2n+2$ is even, together with
$\varphi(2n+3)\leqslant2n+2$. Finally $\sum_{i\leqslant x}\varphi(i)\sim3x^2/\pi^2$,
and summing this once more over $i\leqslant n$ gives the stated asymptotic.
\end{proof}

At length $n$ the reflection machine may therefore be fed any of about
$n^3/\pi^2$ balanced words, of which at most $3n+3$ come from a genuine
occurrence: a proportion $O(n^{-2})$. This is the precise sense in which the
arithmetic language is narrower than the combinatorial one, and it locates
where a sharper form of the results below would have to come from, since
balance is blind to the distinction.

A second consequence of genuineness is available at once, and it removes an
auxiliary choice from the argument of \cref{sec:markoff-application}: the
Farey slope of an occurrence may be used directly as a compatible slope for
its core, with no recourse to the existence statement of
\cref{lem:compatible-slope}.

\begin{corollary}[The Farey slope is compatible]
\label{cor:farey-compatible}
Let $w$ be the core of a genuine reduced occurrence of slope $P/Q$. Then
every nonempty factor $v$ of $w$ satisfies
\begin{equation}
 \bigl\lvert\,\lvert v\rvert_b-(P/Q)\lvert v\rvert\,\bigr\rvert<1.
 \label{eq:farey-compatible}
\end{equation}
\end{corollary}

\begin{proof}
The word $w$ is a factor of $C_{P/Q}$, so by the telescoping computation at
the end of the proof of \cref{prop:genuine-source-dictionary} a factor $v$
of length $r$ has $\lvert v\rvert_b$ equal to $\lfloor rP/Q\rfloor$ or
$\lceil rP/Q\rceil$. Moreover $r\leqslant\lvert w\rvert<Q$ and
$\gcd(P,Q)=1$, so $rP/Q$ is not an integer; each of the two values is
therefore at distance strictly less than one from it.
\end{proof}
